%-------------------------
% classic-serif CV template skeleton
%
% Register with:  /add-template templates/cv/classic-serif/template.tex
% Compile:       lualatex -interaction=nonstopmode -halt-on-error <file>.tex
% Page limit:    exactly 1 page
%
% Layout contract (see TEMPLATE.md for the full style rules):
%   SECTION ORDER IS FIXED: Education -> Work Experience -> Projects ->
%   Achievements and Profile Links -> Technical and Non-Technical Skills.
%   Do not reorder or insert a summary/profile section; Education leads.
%   Section headings are small caps with a rule beneath (via \titleformat).
%   Body font is the default serif — no font package, so it compiles on any
%   TeX install with zero font dependencies.
%-------------------------

\documentclass[letterpaper,10pt]{article}

\usepackage[empty]{fullpage}
\usepackage[dvipsnames]{xcolor}
\usepackage{titlesec}
\usepackage{enumitem}
\usepackage[hidelinks]{hyperref}
\usepackage{fancyhdr}
\usepackage[english]{babel}
\usepackage{tabularx}

\pagestyle{fancy}
\fancyhf{}
\fancyfoot{}
\renewcommand{\headrulewidth}{0pt}
\renewcommand{\footrulewidth}{0pt}

% Margins tuned so a full CV fits exactly one page. If content overflows,
% cut bullet count before loosening these — the page limit is hard.
\addtolength{\oddsidemargin}{-0.55in}
\addtolength{\evensidemargin}{-0.55in}
\addtolength{\textwidth}{1.10in}
\addtolength{\topmargin}{-0.70in}
\addtolength{\textheight}{1.52in}

\urlstyle{same}
\raggedbottom
\raggedright
\setlength{\tabcolsep}{0in}

% Tighten by scaling leading, never by negative \vspace on list items.
% 0.94 keeps the baseline skip (~11.3pt at 10pt) above the descender+ascender
% depth of adjacent lines, so wrapped bullets cannot collide.
\linespread{0.94}

\titleformat{\section}{
  \vspace{-8pt}\scshape\raggedright\large
}{}{0em}{}[\color{black}\titlerule\vspace{-5pt}]

\newcommand{\sepbar}{\hspace{0.45em}\textbar{}\hspace{0.45em}}

% Bold left / right-aligned location + date.
% The p{} column lets a long institution or title wrap rather than overflow.
\newcommand{\resumeSubheading}[4]{
  \vspace{-2pt}\item
  \begin{tabular*}{0.97\textwidth}[t]{p{0.62\textwidth}@{\extracolsep{\fill}}r}
    \textbf{#1} & \textit{\small #2} \\
    \textit{\small #3} & \textit{\small #4} \\
  \end{tabular*}\vspace{-2pt}
}

% Project heading: bold title, tech stack beneath, date right-aligned.
\newcommand{\resumeProjectHeading}[3]{
  \item
  \begin{tabular*}{0.97\textwidth}[t]{p{0.76\textwidth}@{\extracolsep{\fill}}r}
    \small\textbf{#1} & \small\textit{#3} \\
    \textit{\small #2} & \\
  \end{tabular*}\vspace{-2pt}
}

\renewcommand\labelitemii{\textbullet}

\newcommand{\resumeItem}[1]{
  \item\small{#1}%
}

\newcommand{\resumeSubHeadingListStart}{
  \begin{itemize}[leftmargin=0.15in,label={}]
}
\newcommand{\resumeSubHeadingListEnd}{\end{itemize}}
% itemsep is ZERO, not negative. Negative itemsep plus a per-item \vspace pulls a
% wrapped bullet's second line into the next bullet (glyph boxes intersect) while
% the LaTeX log stays clean — an Overfull \hbox will not catch it. Tightness comes
% from \linespread in the preamble instead, which scales leading uniformly.
\newcommand{\resumeItemListStart}{
  \begin{itemize}[leftmargin=0.18in,itemsep=0pt,topsep=0pt,parsep=0pt]
}
\newcommand{\resumeItemListEnd}{\end{itemize}\vspace{-2pt}}

%======================================
\begin{document}

% ── HEADER ───────────────────────────
\begin{center}
  {\Huge \scshape [YOUR_NAME]}\\[3pt]
  \small
  [YOUR_CITY, STATE, COUNTRY]
  \sepbar [YOUR_PHONE]
  \sepbar \href{mailto:[YOUR_EMAIL]}{[YOUR_EMAIL]}
  \sepbar \href{[YOUR_LINKEDIN_URL]}{[YOUR_LINKEDIN_HANDLE]}
  \sepbar \href{[YOUR_GITHUB_URL]}{[YOUR_GITHUB_HANDLE]}
\end{center}

% ── 1. EDUCATION ────────────────────
\section{Education}
\resumeSubHeadingListStart
  \resumeSubheading
    {[YOUR_DEGREE] in [YOUR_FIELD]}
    {[YOUR_CITY, COUNTRY]}
    {[YOUR_INSTITUTION] \textnormal{--- GPA: \textbf{[YOUR_GPA]}}}
    {[YOUR_START] -- [YOUR_END]}
    \resumeItemListStart
      \resumeItem{\textbf{Coursework:} [COURSE 1], [COURSE 2], [COURSE 3], [COURSE 4].}
    \resumeItemListEnd
  \resumeSubheading
    {[YOUR_SECONDARY_QUALIFICATION]}
    {[YOUR_CITY, COUNTRY]}
    {[YOUR_SCHOOL]}
    {[YOUR_START] -- [YOUR_END]}
\resumeSubHeadingListEnd

% ── 2. WORK EXPERIENCE ───────────────
\section{Work Experience}
\resumeSubHeadingListStart
  \resumeSubheading
    {[YOUR_ROLE]}
    {[YOUR_CITY, COUNTRY]}
    {[YOUR_ORGANISATION]}
    {[YOUR_START] -- [YOUR_END]}
    \resumeItemListStart
      \resumeItem{[ACHIEVEMENT WITH A NUMBER: quantify scope, scale, or outcome. Example: Delivered [N] workshops to [M] students with a [P] percent satisfaction rating.]}
      \resumeItem{[SECOND ACHIEVEMENT: lead with the verb, then the mechanism, then the result.]}
    \resumeItemListEnd
\resumeSubHeadingListEnd

% ── 3. PROJECTS ─────────────────────
\section{Projects}
\resumeSubHeadingListStart

  \resumeProjectHeading
    {[PROJECT NAME]}
    {[LANG/FRAMEWORKS] \quad\textbar\quad \href{[PROJECT_URL]}{[PROJECT_HANDLE]}}
    {[START] -- [END]}
    \resumeItemListStart
      \resumeItem{[WHAT IT DOES + the hard part you solved, with a metric. Every bullet carries a number: throughput, latency, scale, coverage, or time saved.]}
      \resumeItem{[DESIGN DETAIL THAT PROVES DEPTH: concurrency model, storage engine, failure handling, or algorithm choice.]}
    \resumeItemListEnd

\resumeSubHeadingListEnd

% ── 4. ACHIEVEMENTS AND PROFILE LINKS
\section{Achievements and Profile Links}
\resumeItemListStart
  \resumeItem{[DISTINCTION WORTH PROOF: a release, an award, a record, a talk, or an open-source contribution. Numbers only if real.]}
  \resumeItem{[SECOND DISTINCTION]}
  \resumeItem{\href{[YOUR_GITHUB_URL]}{[YOUR_GITHUB_HANDLE]} \textbar{} \href{[YOUR_LINKEDIN_URL]}{[YOUR_LINKEDIN_HANDLE]} \textbar{} [CERTIFICATION NAME] ([ISSUER], [MONTH YEAR])}
\resumeItemListEnd

% ── 5. TECHNICAL AND NON-TECHNICAL SKILLS
\section{Technical and Non-Technical Skills}
\begin{itemize}[leftmargin=0.18in,itemsep=-1pt]
  \small
  \resumeItem{\textbf{Programming Languages:} [LANG_1], [LANG_2], [LANG_3].}
  \resumeItem{\textbf{[DOMAIN_2_TITLE]:} [SKILL_1], [SKILL_2], [SKILL_3].}
  \resumeItem{\textbf{[DOMAIN_3_TITLE]:} [SKILL_1], [SKILL_2], [SKILL_3].}
  \resumeItem{\textbf{[DOMAIN_4_TITLE]:} [SKILL_1], [SKILL_2], [SKILL_3].}
  \resumeItem{\textbf{[DOMAIN_5_TITLE]:} [SKILL_1], [SKILL_2], [SKILL_3].}
  \resumeItem{\textbf{Soft Skills:} [SOFT_SKILL_1], [SOFT_SKILL_2].}
\end{itemize}

\end{document}